\begin{frame}{Inheritance — Reusing Common Behavior}

    \begin{center}

        \begin{tikzpicture}[
            font=\small,
            classbox/.style={
                draw=black,
                rounded corners=4pt,
                minimum width=3.8cm,
                minimum height=1cm,
                align=center
            }
        ]

            % Parent class
            \node[
                classbox,
                fill=WorkshopLightBlue
            ] (person) {
                \Large\textbf{Person}\\[3pt]
                name,\ age
            };

            % Child classes
            \node[
                classbox,
                fill=WorkshopGreen,
                below left=1cm and 1.3cm of person
            ] (student) {
                \Large\textbf{Student}\\[3pt]
                student\_id
            };

            \node[
                classbox,
                fill=WorkshopGreen,
                below right=1cm and 1.3cm of person
            ] (teacher) {
                \Large\textbf{Teacher}\\[3pt]
                subject
            };

            % Inheritance arrows
            \draw[
                ->,
                thick
            ]
            (student.north)
            -- (person.south west);

            \draw[
                ->,
                thick
            ]
            (teacher.north)
            -- (person.south east);

        \end{tikzpicture}

    \end{center}

    \vspace{0.2cm}

    \begin{itemize}

        \item Put common attributes and behavior in a parent class

        \item Child classes inherit that common functionality

        \item Child classes can add their own attributes and behavior

    \end{itemize}

    \vfill

    \begin{center}
        \textbf{Person}
        $\rightarrow$
        \textbf{Student / Teacher}
    \end{center}

\end{frame}